% Nouveau EGA V (cote 25): archive cote 25, site page 129 (PDF page 130).
% typescript. Draft compared: none.
% AI-transcribed by Claude Opus 5.5 (claude-opus-5-5, adaptive thinking, effort high), Message Batches API, batch ega25b (msgbatch_01A7n5YuXubUPNK1zsepi8RG), 2026-11-13; not yet checked by eye.
\scanpage{25/129}
\mspage{\hand{32}}
% NOTE: A pencil oblique stroke follows the upper page number.

précédent, et de la détermination des automorphismes de $\overset{\hand{\vee}}{\mathrm{P}}(E) = P(\overset{\hand{\vee}}{E})$ faite
au n°1 , qui implique qu'un automorphisme d'un fibré projectif sur un S
(relatif à un Module loc libre de type fini)  est connu quand on connait
son effet sur une famille de sections formant une ``base projective''
(i.e. formant une base projective en chaque  fibre).

Esquissons maintenant le cas général: S quelconque.  Bien entendu,
quitte à faire un changement de base sur S, on est ramené dans a) à prouve
que deu sections  de $\overset{\hand{\vee}}{\mathrm{P}}$ sur S  qu\uline{i} définissent le même diviseur sur X sont
identiques, et dans b) que deux automorphismes  de $\overset{\hand{\vee}}{\mathrm{P}}$  qui sont tels que
$D\circ u = D\circ v$  sont identiques. On peut supposer  S affine, Dans le cas b) où on n'a pas supposé expressé%
ment  g de présentation finie, mais g surjectif, on se ramène aussitôt
(grâce au fait que g est ouvert) au cas où X est également affine, donc
de présentation finie sur S.  Par le procédé bréveté, on peut maintenant
se réduire au cas S \uline{noethérien} .  Or pour un préschéma de base Sloc noethérien,
et un morphisme de foncteurs  $F \overset{\hand{h}}{\hand{\longrightarrow}} G$ sur S  (F et G des foncteurs
$(\mathrm{Sch})_{/S}^{\circ} \hand{\longrightarrow} (\mathrm{Ens})$ ) ,  on a \unsure{un} critère très généra\unsure{l} (qui sera explicité au Chap V), permettant d'affir-%
mer que si  pour tout  $s \hand{\in} S$, le morphisme correspondant $F_s \hand{\longrightarrow} G_s$ est
un monomorphisme, alors $F \hand{\longrightarrow} G$ est un monomorphisme  (NB on pose $F_s{=}F\times_S\mathrm{Spec}$
$(k(s))$, et de même pour $G_s$ ),  moyennant des hypothèses simples sur F et G
(vérifiées par exemple si F et G sont tous les deux représentables par des
préschémas loc de type fini sur S; mais dans le cas qui nous occupe, seul
le premier des deux foncteurs est représentable à priori).  Nous allons
esquisser le raisonnement du Chap V  dans le deux cas particuliers qui nous
intéressent ici.  On a deux sections u, v (de $\overset{\hand{\vee}}{\mathrm{P}}$ , resp. du groupe projectif
\uline{GP} ($\overset{\hand{\vee}}{E}$) ) d'un \hand{pré}schéma de type fini F sur S dont on veut montrer qu'elles
sont égales, pour ceci manifestement il suffit de prouver qu'elles sont
égales après tout changement de base  $\mathrm{Spec}(\underline{O}_{S,s}/\underline{m}_s^{\hand{n+1}}$ ) $\hand{\longrightarrow}$ \hand{$\mathcal{S}$},

\mspage{\hand{129}}
